\section{Discussion and Open Problems}
\label{sec:discussion}

We asked whether QIPMs beat classical IPMs on sparse problems.  Instead of a
single yes or no, the results separate three levels, each with its own
obstructions and ways around them, that informal complexity claims often
combine: the geometry of the central path, the model of access to the
resulting linear-algebra objects, and the interface through which one
iteration produces the next.

\subsection{What the paper establishes}
\label{sec:discussion-verdict}

\paragraph{Geometry.}
For barrier-Newton methods, part of the condition-number growth is forced
by the optimization problem.  \Cref{thm:sublevel-chord-law} bounds the
condition number of the reduced barrier Hessian from below at a fixed
objective gap, for every self-concordant barrier, and
\cref{thm:canonical-achievability} shows that the log and log-det barriers
attain the resulting order up to instance-dependent constants.
\Cref{thm:width-spectrum-rigidity} extends this to a geometric description
of the whole spectrum.  In the stated minimax sense, changing the barrier
cannot in general remove the lower bound given by sublevel chords.

This is not the folklore claim that every central path has
\(\kappa=\Theta(\mu^{-2})\); the rate depends on the geometry of the
sublevel sets.  A positive-dimensional optimal face gives growth of order
\(g^{-2}\) in the objective gap \(g\); a curved unique optimum can give
growth of order \(g^{-1}\); at a unique LP vertex the condition number can
stay bounded; and an attained singularity degree can give fractional
exponents (\cref{cor:conditioning-regimes,cor:fractional-conditioning-laws}).
Nor does a large condition number alone determine solver cost: Krylov
iteration counts can be polylogarithmic in an arbitrarily large ratio
between two eigenvalue clusters
(\cref{prop:two-cluster-hessian,prop:two-cluster-cg}), and the Newton
right-hand side on the path can avoid the small eigenvalues
(\cref{prop:newton-rhs-alignment,cor:top-window-solve}).

\paragraph{Access models.}
The quantum cost depends on how the matrix is supplied, not only on its
spectrum.  The clustered comparison in \cref{thm:quantum-cluster-obstruction}
separates plain normalized block-encoding access from access to a factor
\(G\) with \(H=G^\dagger G\).  With plain access, the known worst-case query cost
\(\Thetatilde(\kappa)\) remains for two-point spectra in sufficiently large
dimension.  Norm-tight factor access admits a polynomial of square-root cost
but an end-to-end solve only under overlap assumptions; a generic
factor-access QLS solver of square-root cost would contradict the
plain-access lower bound, because the square-root factor of that hard
instance is easy to query.  By the exact coupling and
spectral-dimension laws
(\cref{thm:dls-coupling,thm:spectral-dimension-laws}), the distribution of the
right-hand side over the spectrum decides whether truncation or filtering
can improve on \(\kappa\).

The same holds for preconditioning.  A whitening factor may trade condition
number for a larger normalization; a congruence may trade the transformed
condition number for sensitivity when the physical solution is recovered; and
a perfect preconditioner that depends on the input may move the hard
information into setup or RHS preparation
(\cref{thm:prec-whitening-frontier,thm:prec-congruence-frontier,thm:prec-factor-rhs-hardness,thm:prec-parity-ledger}).
These results bound products of costs or the total cost of all phases.
They show not that preconditioning is useless, but that a small transformed
condition number is not by itself a complexity bound.

\paragraph{Interfaces.}
In our hard families the main obstruction is often loading or recovery
rather than inversion.  The prefix LP has a reduced Hessian of condition
number one along its central path, yet preparing a robust primal direction
in the original coordinates needs linearly many raw queries
(\cref{thm:lp-prefix-capstone}).  The SDP families distinguish a gap in the
optimal value (\cref{thm:sdp-intrinsic}), parity-hard primal and
complete-KKT solution states (\cref{thm:sdp-newton-state}), and a public Schur
solve for which loading the original coordinates obeys an accuracy--query
tradeoff (\cref{thm:sdp-schur,thm:sdp-accuracy-frontier}).  The canonical KKT
block encoding is exact at one raw query per call, yet preparing the root,
global, or complete-KKT state still needs linearly many calls to it, and
first constructing a block encoding of the tangent projector needs linearly
many raw queries (\cref{sec:sdp-intrinsic-access}).  Nor is a normalized state a
reusable coordinate oracle: producing the changing diagonal of the next Newton system
obeys the state-to-diagonal and state-to-next-solve product bounds in
\cref{thm:prec-state-diagonal,thm:prec-state-next-solve}.

In the condensation results, the winner mass controls both the condition
number of the reduced Hessian and the trace
distance from a public state with no winner
(\cref{thm:cond-mass-conditioning,thm:cond-visibility}).  The bounded-incidence
holonomy SDP has matching query bounds for the value and for a compact
output that names the winner (\cref{thm:cond-unique-or,thm:wta-search-upper}),
while the cost of preparing the central state changes across the phases of
condensation (\cref{sec:wta-central-state,eq:wta-preparation-phases}).  The
decoder that verifies sampled labels turns a trace-close mixed output into a
public decision procedure with one-sided error, and parity collapse gives
the query cost as a function of the allowed trace-distance error
(\cref{thm:sc-m1,thm:sc-small-success,cor:sc-s2}).  The lower bound holds for
a specified accuracy of the output state; a constant primal objective gap
alone would not imply it.

In short, conditioning is governed by sublevel geometry, quantum
performance by the access actually supplied, and end-to-end performance by
loading, recovery, updates, checks, and output between iterations.  Sparsity
can help at each level, but no single notion of sparsity controls all
three.

\subsection{What the paper does not establish}
\label{sec:discussion-nonclaims}

The paper proves neither an unconditional sparse-QIPM advantage nor an
unconditional lower bound.  The Apers--Gribling result for tall LPs remains
the only proved asymptotic quantum--classical separation, in row queries,
among the QIPM results considered here
\cite{apers2026-quantum-speedups-linear-programming-interior}; our results do
not weaken it.  Conversely, no parity or holonomy family rules out a method
with a different oracle, a weaker but still useful output, an
optimization primitive free of the condition number, or a structure-aware
preconditioner outside the stated class.

The barrier theorem concerns the reduced Hessian in an orthonormal basis of
the tangent space, not automatically an augmented KKT system, primal--dual
scaling, or a Hamiltonian central-path method.  The clustered QLS theorem
depends on the access model and has conditions on the dimension and on
overlaps.  The LP and SDP loading theorems count raw
coefficient queries; a supplied parity, prefix, tangent-projector, or
preconditioned-product oracle is a stronger input model.  The online direct
sum requires fresh independent input in each round, whereas a fixed instance
can be cached (\cref{thm:prec-online-direct-sum,sec:preconditioning-lazy}).
The physical light-cone bound concerns data and output placed locally with
bounded occupancy per site, not all-to-all query access
(\cref{thm:structural-light-cone}).

Likewise, the positive modules are components, not a complete algorithm.
Projective reuse of a rescaled Newton direction, as analyzed here, needs
finite realized variation and a fixed residual tolerance; equilibration of
the active range and face repair need known strict-complementarity tail margins and spectral gaps;
compressed dual output needs a retained dimension that is actually smaller
and must count the cost of its coordinate compiler; and the block-angular
module needs a projected barrier with a small parameter and a local
representation kept between iterations (\cref{sec:positive-synthesis}).  The
counterexamples in
\cref{sec:positive-lazy,sec:positive-face-repair,sec:positive-checkpoint-span}
show that the corresponding module fails when one of these conditions is
removed.

We propose the five accounting rules of \cref{sec:positive-modules} as a
standard for future positive claims, and add the comparison rules used
throughout the paper: match the residual
and precision requirements, count setup and reuse costs on both sides, and
compare with the best classical method for the same structure.  A query
advantage is not automatically a gate or time advantage, and a comparison of
two upper bounds is not a separation.

\subsection{Collected open problems}
\label{sec:discussion-open}

We collect every open-problem environment in
\cref{sec:conditioning,sec:spectra,sec:state-conversion,sec:lp-hardness,sec:sdp-hardness}
and the further open questions raised in the technical discussion, grouped
by the level at which new information is needed.

\subsubsection{Geometry and conditioning}

\begin{enumerate}[leftmargin=*,label=G\arabic*.]
 \item \textbf{Degenerate-vertex constants.}  Determine the limiting spectrum,
 not only its boundedness, of the canonical reduced log-barrier Hessian at a
 unique degenerate LP vertex from the geometry of the active set
 (\cref{open:degenerate-vertex-constants}).
 \item \textbf{Facial characterization of chord growth.}  Characterize the
 exact growth \(D(g)\) of objective-sublevel chords of an SDP from its facial
 data, and decide when the singularity-degree exponent is attained by the
 sublevel diameters
 (\cref{open:facial-chord-growth}).
 \item Relate the neighborhood \eqref{eq:cond-robust-residual} used by robust
 condensation, which bounds the blockwise stationarity residual in the
 barrier norm by \(\eta\) and the trace residual by \(\zeta\), to the
 algorithm-specific inexact Newton residuals of
 \cref{def:inexact-newton-residual}, whose forcing tolerance is a separate
 \(\eta\) (\cref{sec:condensation-discussion}).
\end{enumerate}

\subsubsection{Quantum access and state conversion}

\begin{enumerate}[leftmargin=*,label=A\arabic*.]
 \item \textbf{Exact subnormalization shift.}  Decide whether plain access to
 \(H\) permits a block encoding of \(I-H\) with unit subnormalization without
 paying the apparent cost of uniform amplification at the shifted spectral
 edge (\cref{open:exact-subnormalization-shift}).  This is the
 promise-sensitive form of the normalization question of Orsucci and Dunjko
 \cite{orsucci2021-on-solving-classes-of-positive}.
 \item \textbf{General inner predicate.}  Prove or refute the
 \(\sqrt\delta\,Q(f)\sqrt{G/k}\) small-success composition lower bound for a
 general inner predicate \(f\), beyond the proof by parity collapse
 (\cref{open:sc-general-composition}).
 \item Determine the minimax output radius and query curve when the public
 internal states of the losing blocks differ by \(\eta_{\rm int}\), including
 the regime below condensation (\cref{sec:state-conversion-open}).
 \item Resolve the state-conversion range \(\delta=o(k/G)\), where a public
 output with uniform labels may already satisfy part of the output requirement
 and the winner mass alone does not determine the boundary
 (\cref{sec:state-conversion-open}).
 \item \textbf{Intrinsically nonabelian word hardness.}  Prove a lower bound
 on local-symbol queries for distinguishing conjugacy classes under a
 promise on words whose letters do not commute, and transfer it to output of
 the conjugacy class in the group-holonomy SDP, instead of restricting to a
 cyclic subgroup
 (\cref{open:sdp-nonabelian}).
\end{enumerate}

\subsubsection{Hard families and classes of gadgets}

\begin{enumerate}[leftmargin=*,label=F\arabic*.]
 \item \textbf{Bounded-scale full-KKT parity amplifier.}  Construct, or rule
 out, a standard-form LP of linear size with bounded row and column degree;
 bounded coefficients, each depending on at most one input bit; public Newton
 right-hand sides that are simple to prepare; primal and slack starting
 values bounded above and below by positive constants; a complete Newton
 direction bounded coordinatewise; and a parity-hard primal component with
 constant mass in the state of that complete direction
 (\cref{open:lp-bounded-full-kkt}).
 \item Determine whether an arbitrary sparse equality LP with bounded range
 can support a nonlinear decoder of the normalized state with a fixed global
 residual tolerance without introducing nonlocal coefficient access
 (\cref{sec:frontier-synthesis}).
 \item Determine whether a nonorthogonal mixed-mode gadget can solve the same
 bounded-scale complete-direction problem while counting the cost of every
 input-dependent source and recovery map
 (\cref{sec:lp-beyond-pairs,sec:frontier-synthesis}).
\end{enumerate}

\subsubsection{Positive algorithmic results}

\begin{enumerate}[leftmargin=*,label=P\arabic*.]
 \item Exhibit one explicit family and access model that satisfy all six
 requirements of \cref{sec:positive-synthesis} simultaneously: stable direct
 normalization and inverse, stable forcing and RHS preparation, finite
 projective variation or a barrier with a small parameter, cheap static
 construction, useful compressed output, and a classical lower bound that
 holds even with randomized preprocessing.
 \item Find a structured access model that supplies a quantum operator with
 constant normalization without also supplying a classical sampling
 distribution with bounded leverage, and realize it on an explicit sparse LP
 family (\cref{sec:positive-checkpoint-span}).
 \item Construct a projected block-angular barrier with parameter
 \(\Otilde(k+\poly(p))\) whose local center and derivatives can be recomputed
 from one scenario block and the explicit border
 (\cref{sec:positive-block-angular}).
\end{enumerate}

The list omits directions that a theorem here already rules out: turning a
bounded-degree local-linear parity gadget into an
expander does not escape the mixture and resistance bounds, and ordinary
residual control in short steps does not imply finite projective variation
(\cref{sec:frontier-synthesis,sec:positive-lazy}).  A new proposal must drop a
premise of these theorems, not relabel the same construction.

\subsection{Outlook}
\label{sec:discussion-outlook}

The most useful next result would be small and complete: one structured
family, one explicit access model, one useful
output requirement, an accounting of all costs along the path, and a matching
classical lower bound or algorithmic comparison.  Problem~P1 is deliberately
hard because each module alone already has a classical or quantum
obstacle.  Stable spectral geometry is not enough if RHS preparation is
hard; cheap coherent access is not enough if a classical algorithm can
sample the same averages; and compressed output is not enough if the next
iterate must be written out in full.

Among nearer-term targets, Problem~A1 would show whether the separation
between plain and factor access for clustered QLS is an artifact of the
current construction or a property of the access models; Problem~A2 would
make the state-conversion curve a reusable composition principle rather than
a result about parity; Problem~A5 would separate matrix-valued holonomy from
its abelian restriction; and Problem~P3 would connect the border module to a
candidate end-to-end QIPM.

The methodological lesson is that a future sparse-QIPM advantage cannot be
claimed from one favorable parameter: it must hold for the geometry, the
access model, and the output interface at once.  This paper proves bounds
and counterexamples for each of these, with conditional modules where one
can be made cheap.  Solving Problem~P1 would give the end-to-end answer that
this paper does not claim.
